% ch04  §5.2–§5.4 / §6.1–§6.3（原文 1648–2027 行）

\subsection{Hecke 作用}

设 $V\in \Rep(\dG^{I})$。

\ssszh{}
对任意 $x\in \Sig$，$G(F_{x})$ 上具紧支集、双 $K_{x}$-不变的函数构成的向量空间 $\bH_{x}:=\Qlbar[K_{x}\bs G(F_{x})/K_{x}]$ 在卷积下是 $\Qlbar$-代数，单位元为 $\one_{K_{x}}$。有环同态
\begin{equation*}
\bH_{x}\to \End(\cH^{I}(V))
\end{equation*}
定义如下。

由推论 \ref{c:same HV}，只需在 $\cH^{(I_{1},\cdots, I_{r})}(V)$ 上构造 $\cH_{x}$ 的作用\zy{原文此处作 $\cH_{x}$；据同段前文，作用在 $\cH^{I}(V)$ 上的应是 Hecke 代数 $\bH_{x}$，$\cH$ 与 $\bH$ 疑为笔误。}，其中各 $I_{j}$ 都是单点集。又可将情形化归为 $V$ 不可约，此时 $V=\boxtimes_{i\in I}V_{i}$，其中 $V_{i}$ 是 $\dG$ 的不可约表示，最高权为 $\l_{i}$。令 $\un\l=(\l_{i})_{i\in I}$。对 $[g_{x}]\in K_{x}\bs G(F_{x})/K_{x}$，回忆 Hecke 对应 \eqref{Hk x Sig}，其中 $\oll{c}$ 与 $\orr{c}$ 由引理 \ref{l:change level Sig} 都是有限 étale 的。

类似地可以构造相对位置映射
\begin{equation*}
\ev_{\Sht[g_{x}]}^{(I_{1},\cdots, I_{r})}: \Sht^{(I_{1},\cdots, I_{r})}_{G,\bK}[g_{x}]\to (\QGr_{G}/\Aut(D))^{I}
\end{equation*}
它与映射 $\oll{c}$、$\orr{c}$ 以及 \eqref{ev Sht singleton} 中 $\Sht^{(I_{1},\cdots, I_{r})}_{G,\bK}$ 的相对位置映射相容。于是得到典范同构
\begin{equation}\label{corr IC}
\oll{c}^{*}\IC_{\Sht}^{I}(V)\isom \orr{c}^{*}\IC_{\Sht}^{I}(V)\cong\orr{c}^{!}\IC_{\Sht}^{I}(V).
\end{equation}
$\one_{K_{x}g_{x}K_{x}}$ 在 $\cH^{I}(V)$ 上的作用定义为如下复合：
\begin{equation*}
\begin{aligned}
\cH^{I}(V)=\pi^{I}_{!}\IC_{\Sht}^{I}(V)&\to \pi^{I}_{!}\oll{c}_{!}\oll{c}^{*}\IC_{\Sht}^{I}(V)\\
&\xr{\eqref{corr IC}} \pi^{I}_{!}\orr{c}_{!}\orr{c}^{!}\IC_{\Sht}^{I}(V)\to \pi^{I}_{!}\IC_{\Sht}^{I}(V)=\cH^{I}(V)
\end{aligned}
\end{equation*}
这里用到了伴随的单位与余单位：$\id \to \oll{c}_{*}\oll{c}^{*}\cong \oll{c}_{!}\oll{c}^{*}$（因为 $\oll{c}$ 是 proper 的），以及 $\orr{c}_{!}\orr{c}^{!}\to\id$。

\ssszh{球 Hecke 作用}
对 $x\in |U|=|X|-\Sig$，类似的构造给出环同态
\begin{equation*}
\bH_{x}:=\Qlbar[G(\cO_{x})\bs G(F_{x})/G(\cO_{x})]\to \End(\cH^{I}(V)|_{(U-\{x\})^{I}}).
\end{equation*}

\begin{rmk} 对 $x\in |U|$，$\bH_{x}$ 在 $\cH^{I}(V)|_{(U-\{x\})^{I}}$ 上的作用已被 V. Lafforgue 延拓到整个 $\cH^{I}(V)$ 上：办法是把它看成游走算子的一个特例。
\end{rmk}

设 $\Sig'\subset |U|$ 是任意有限个位点构成的集合。把上面定义的各 Hecke 作用合起来，得到环同态
\begin{equation*}
\bH_{\Sig'}=\ot_{x\in \Sig'}\bH_{x}\to \End(\cH^{I}(V)|_{(U-\Sig')^{I}}).
\end{equation*}
取 $\Sig'$ 递增时的余极限，就得到整体 Hecke 代数在 $\cH^{I}(V)$ 限制到一般点 $\y^{I}\to X^{I}$ 之上的作用
\begin{equation*}
\bH_{G,\bK}:=\ot'_{x\in |X|}\bH_{x}\to \End(\cH^{I}(V)|_{\y^{I}}).
\end{equation*}

\subsection{因子化结构}

\ssszh{预期的形态}
设 $V=\boxtimes_{i\in I}V_{i}$，其中 $V_{i}\in \Rep(\dG)$。则\emph{预期} $\cH^{I}(V)$ 满足下述因子化性质。设 $\pi$ 是 $G(\AA)$ 的尖点自守表示。令 $\pi(\bK)\subset \pi$ 为在 $\prod_{x\notin\Sig}G(\cO_{x})\times \prod_{x\in \Sig}K_{x}$ 下不变的那些向量构成的子空间，其中 $K_{x}=\bK_{x}(k_{x})$。注意 $\pi(\bK)$ 是 $\bH_{G,\bK}$-模。由 V. Lafforgue 的工作 \cite{VLaff}，可以对 $\pi$ 附上一个半单 Langlands 参数 $\r_{\pi}$，它是 $U$ 上的 $\dG$-局部系统。利用 V. Lafforgue 的游走算子，可以定义 $H^{i}\cH^{I}(V)$（即 $\cH^{I}(V)$ 的第 $i$ 个上同调层，$i\in \ZZ$）的 $\pi$-等型子层，记作 $H^{i}\cH^{I}(V)[\pi]$。作为初步近似，可以把 $H^{i}\cH^{I}(V)[\pi]$ 定义为 $\Hom_{\bH_{G,\bK}}(\pi(\bK), H^{i}\cH^{I}(V))$（它可能比真正的 $H^{i}\cH^{I}(V)[\pi]$ 更大；微妙之处在于 $\bH_{G,\bK}$-模 $\pi(\bK)$ 未必能决定 $\r_{\pi}$ 的同构类）。

我们作一个简化假设：$\Aut(\r_{\pi})=Z\dG$。例如当 $G=\GL_{n}$ 且 $\pi$ 尖点时（此时 $\r_{\pi}$ 是不可约的秩 $n$ 局部系统），已知确实如此。于是预期
\begin{equation*}
H^{0}\cH^{I}(V)[\pi]\stackrel{?}{\simeq} M_{\pi}\ot \boxtimes_{i\in I}(V_{i,\r_{\pi}}),
\end{equation*}
其中 $V_{i,\r_{\pi}}$ 是由 $\r_{\pi}$ 与 $\dG$-表示 $V_{i}$ 诱导出的 $U$ 上的局部系统，$M_{\pi}$ 是依赖于 $\pi$ 的有限维向量空间。还预期此时 $H^{i}\cH^{I}(V)[\pi]=0$ 对 $i\ne0$ 成立。更精确的陈述是数域上自守形式的 Arthur 重数公式的类比，见 Xinwen Zhu 的文章 \cite[\S4.7]{Zhu}。

上述预期特别地暗示：在 $\cH^{I}(V)$ 的几何一般茎上，有 $I$ 份基本群 $\pi_{1}(U,\ov \y)$ 的相互交换的作用。这一事实由 Cong Xue \cite{X} 证明，推广了 V. Lafforgue \cite{VLaff} 的结果。他们的证明依赖部分 Frobenius（将在 \S\ref{sss:pF} 中讨论）与 Drinfeld 引理。

\ssszh{在部分对角线上的行为}
下面把 $\cH^{I}(V)$ 限制到部分对角线上时的行为讲清楚。

有限集之间的映射 $\th: I\to J$ 诱导群同态 $\D_{\dG, \th}: \dG^{J}\to \dG^{I}$，从而诱导表示上的限制函子
\begin{equation*}
\Res_{\th}: \Rep(\dG^{I})\to \Rep(\dG^{J}).
\end{equation*}

\begin{prop}
设 $\th: I\surj J$ 是有限集之间的满射，相应地有部分对角映射 $\D_{\th,U}: U^{J}\incl U^{I}$。则有典范同构
\begin{equation*}
\D_{\th,U}^{*}\cH^{I}(V)\isom \cH^{J}(\Res_{\th}V).
\end{equation*}
此外，该同构对 $V\in \Rep(\dG^{I})$ 具有函子性，并与满射的复合 $I\surj J\surj K$ 相容。
\end{prop}
\begin{proof}[证明概要]
由 proper 基变换，只需对 $\Hk^{J}_{G,\bK}$ 上的 Satake 层验证相应的陈述：
\begin{equation*}
\IC^{I}_{\Hk}(V)|_{\Hk^{J}_{G,\bK}}\cong \IC^{J}_{\Hk}(\Res_{\th}(V)).
\end{equation*}
这里用同构 \eqref{Hk diagonal} 把 $\Hk^{J}_{G,\bK}$ 看作 $\Hk^{I}_{G,\bK}$ 的子叠（经由 $\D_{\th, U}: U^{J}\incl U^{I}$ 拉回）。取双射 $\s: J\isom \{1,\cdots, r\}$，并对 $1\le j\le r'$ 令 $I_{j}=\th^{-1}(\s'^{-1}(j))$。于是有笛卡尔方块
\begin{equation*}
\begin{tikzcd}
\Hk^{(\{1\},\cdots, \{r\})}_{G,\bK} \arrow[r,"\y"] \arrow[d] & \Hk^{(I_{1},\cdots, I_{r})}_{G,\bK} \arrow[d] \\
\Hk^{J}_{G,\bK} \arrow[r] & \Hk^{I}_{G,\bK}
\end{tikzcd}
\end{equation*}
由 proper 基变换与引理 \ref{l:IC Hk refine}，只需证明
\begin{equation}\label{IC fact b}
\y^{*}\IC^{(I_{1},\cdots, I_{r})}_{\Hk}(V)\cong \IC^{(\{1\},\cdots ,\{r\})}_{\Hk}(\Res_{\th}V).
\end{equation}

下面处理 $J$ 本身是单点集的情形。一般情形不再需要新的想法。于是设 $J=\{1\}$ 为单点集。取双射 $\s: I\isom \{1,\cdots, s\}$，并定义
\begin{equation*}
\Hk^{[s]}_{G,\bK}=\Hk^{(\{1\},\cdots,\{s\})}_{G,\bK}\times_{U^{s}}\D(U).
\end{equation*}
于是 $\Hk^{[s]}_{G,\bK}$ 分类带 $\bK$-水平结构的 $G$-主丛的修正 $\cE_{0}\dxr\cE_{1}\dxr\cdots\dxr\cE_{s}$，且这些修正都落在同一个点 $x\in U$ 上。令 $\QGr^{[s]}_{G}$ 是这样的模叠：其 $R$-点是圆盘 $D_{R}$ 上的 $G$-主丛 $(\cE_{0},\cdots, \cE_{s})$，并带有同构 $\a_{i}: \cE_{i-1}|_{D_{R}^{\times}}\isom\cE_{i}|_{D^{\times}_{R}}$（$1\le i\le s$）。于是有求值映射 $\ev^{[s]}: \Hk^{[s]}_{G,\bK}\to \QGr^{[s]}_{G}/\Aut(D)$，以及一个三个方块全为笛卡尔的图表
\begin{equation*}
\begin{tikzcd}
 & \QGr^{[s]}/\Aut(D) \arrow[rr,"\d"] \arrow[dd,"\g"' pos=0.22] & & (\QGr/\Aut(D))^{s} \\
\Hk^{[s]}_{G,\bK} \arrow[rr,"\wt\y" pos=0.62] \arrow[dd] \arrow[ur,"\ev^{[s]}"] & & \Hk^{(\{1\},\cdots,\{s\})}_{G,\bK} \arrow[ur,"\ev^{(\{1\},\cdots,\{s\})}"] \arrow[dd] & \\
 & \QGr/\Aut(D) & & \\
\Hk^{\{1\}}_{G,\bK} \arrow[ur,"\ev^{\{1\}}"] \arrow[rr,"\y"] & & \Hk^{I}_{G,\bK} &
\end{tikzcd}
\end{equation*}
映射 $\d$ 由 $\QGr^{[s]}\to (\QGr)^{s}$ 诱导，后者把 $((\cE_{i})_{0\le i\le r},(\a_{i})_{1\le i\le r})$ 送到 $(\cE_{i-1},\cE_{i}, \a_{i})_{1\le i\le s}$；映射 $\g$ 由 $\QGr^{[s]}\to \QGr$ 诱导，后者把 $((\cE_{i})_{0\le i\le r},(\a_{i})_{1\le i\le r})$ 送到 $(\cE_{0},\cE_{s}, \a_{r}\c\cdots\c\a_{1})$。

再由 proper 基变换，\eqref{IC fact b} 化归为同构
\begin{equation*}
\g_{!}\d^{*}\Sat_{\{1,\cdots , s\}}(V)\cong \Sat(\Res^{\dG^{s}}_{\dG} V).
\end{equation*}
由线性性，只需对 $V=\boxtimes_{i=1}^{s} V_{i}$ 证明；此时左端是张量范畴 $\Perv_{\Aut(D)}(\QGr)$ 中的卷积 $\Sat(V_{1})\star\cdots\star\Sat(V_{s})$，右端是 $\Sat(V_{1}\ot\cdots\ot V_{s})$。所求同构来自 Satake 等价的幺半结构。
\end{proof}

\subsection{部分 Frobenius}\label{sss:pF}
设 $I=I_{1}\coprod\cdots\coprod I_{r}$ 是分解为\emph{单点集}的分解。这给出双射 $\s: I\isom \{1,\cdots, r\}$，使 $I_{j}=\{\s^{-1}(j)\}$。回忆 \eqref{pl Fr} 中定义的部分 Frobenius 映射 $\Fr^{\pl}$。
我们有交换图表
\begin{equation}\label{pl Frob ev}
\begin{tikzcd}
\Sht^{(I_{1},\cdots, I_{r})}_{G,\bK} \arrow[r,"\ev^{(I_{1},\cdots, I_{r})}"] \arrow[d,"\Fr^{\pl}"'] & (\QGr/\Aut(D))^{I} \arrow[d,"{(\Fr, \id,\cdots, \id)}"] \\
\Sht^{(I_{2},\cdots, I_{r}, I_{1})}_{G,\bK} \arrow[r,"\ev^{(I_{2},\cdots, I_{1})}"] & (\QGr/\Aut(D))^{I}
\end{tikzcd}
\end{equation}
这里右侧的竖直映射在第一个因子 $I_{1}$ 上是 Frobenius，在其余因子上是恒等映射。

对最高权为 $\l$ 的不可约表示 $V_{\l}\in \Rep(\dG)$，在 $\Sat(V_{\l})$ 上固定一个 Weil 结构，使其在 $\Gr_{G,\l}$ 上的限制为 $\Qlbar[\j{2\r, \l}](\j{\r,\l})$（若 $\j{\r,\l}$ 是半整数，这里涉及 $q^{1/2}$ 的选取）。用加性把这些 Weil 结构延拓到所有 $V\in \Rep(\dG)$ 的 $\Sat(V)$ 上。更一般地，对任意有限集 $I$ 与 $V\in \Rep(\dG^{I})$，这些 Weil 结构延拓为 $\Sat_{I}(V)\in \Perv((\QGr/\Aut(D))^{I})$ 上的 Weil 结构，并与外张量积相容。

现在设 $V=\boxtimes_{i\in I}V_{i}\in \Rep(\dG^{I})$。对 $1\le j\le r$，把 $V_{\s^{-1}(j)}$ 也记作 $V_{j}$。由图表 \eqref{pl Frob ev} 并利用上面固定的 $\Sat(V_{1})$ 上的 Weil 结构，得到同构
\begin{equation}\label{pl Fr pullback IC}
\IC^{(I_{1},\cdots, I_{r})}_{\Sht}(V)\cong \Fr^{\pl*}\IC^{(I_{2},\cdots, I_{1})}_{\Sht}(V).
\end{equation}
另一方面，对拓扑笛卡尔图表 \eqref{pl Fr UI} 用 proper 基变换，给出同构
\begin{equation*}
\Fr_{I_{1}}^{*}\cH^{I}(V)= \Fr_{I_{1}}^{*}\pi^{(I_{2},\cdots, I_{1})}_{!}\IC^{(I_{2},\cdots, I_{1})}_{\Sht}(V)\isom  \pi^{(I_{1},\cdots, I_{r})}_{!}\Fr^{\pl*}\IC^{(I_{2},\cdots, I_{1})}_{\Sht}(V)
\end{equation*}
结合 \eqref{pl Fr pullback IC} 便得到同构
\begin{equation*}
\phi_{I_{1}}: \Fr_{I_{1}}^{*}\cH^{I}(V)\cong \pi^{(I_{1},\cdots, I_{r})}_{!}\Fr^{\pl*}\IC^{(I_{2},\cdots, I_{1})}_{\Sht}(V)\cong \pi^{(I_{1},\cdots, I_{r})}_{!}\IC^{(I_{1},\cdots, I_{r})}_{\Sht}(V)=\cH^{I}(V).
\end{equation*}
该同构与 $I-I_{1}$ 分解为单点集的方式无关。此外，这类同构按加性延拓到所有 $V\in \Rep(\dG^{I})$。

由于可以取双射 $\s$ 使 $I_{1}$ 成为任意给定的元素 $i\in I$，故对一切 $i\in I$ 与 $V\in\Rep(\dG^{I})$ 有典范同构
\begin{equation*}
\phi_{i}:  \Fr_{i}^{*}\cH^{I}(V)\isom  \cH^{I}(V).
\end{equation*}
这里 $\Fr_{i}: U^{I}\to U^{I}$ 在第 $i$ 个因子上是 Frobenius，在其余因子上是恒等映射。

下面的事实由定义容易验证。
\begin{prop}
同构族 $(\phi_{i})_{i\in I}$ 彼此在典范同构的意义下交换，它们的复合是 $\cH^{I}(V)$ 上的典范 Weil 结构，即由 $\Sat_{I}(V)$ 上的 Weil 结构诱导出的同构 $\Fr^{*}_{U^{I}}\cH^{I}(V)\isom  \cH^{I}(V)$。
\end{prop}

特别地，对任意子集 $J\subset I$ 有典范同构 $\phi_{J}: \Fr_{J}^{*}\cH^{I}(V)\isom  \cH^{I}(V)$，其中 $\Fr_{J}: U^{I}\to U^{I}$ 在 $U^{J}$ 上是 Frobenius，在其余因子上是恒等映射。它定义为各 $\phi_{i}$（$i\in J$）按任意次序的复合。

\section{进阶课题}\label{s:adv}
本节概述 shtuka 模空间的若干方面，这些方面值得进一步研究。

\subsection{紧化}
映射 $\pi^{(I_{1},\cdots, I_{r}),\le \un\l}_{G,\bK}: \Sht^{(I_{1},\cdots, I_{r}),\le \un\l}_{G,\bK}\to U^{I}$ 不是 proper 的，这给研究它的直接像复形、尤其是其通用纤维的上同调带来困难。例如，如果它是 proper 的，复形 $\cH^{I}(V)$ 就会自动是纯的，这将对函数域上自守形式的 Ramanujan–Petersson 猜想产生影响。

由于 $\pi^{(I_{1},\cdots, I_{r}),\le \un\l}_{G,\bK}$ 不是有限型的，要作紧化，必须先限制到 Harder–Narasimhan 截断 ${}^{\le \mu}\Sht^{(I_{1},\cdots, I_{r}),\le \un\l}_{G,\bK}$，如 \S\ref{sss:HN} 那样。

在 \cite{Dr-comp} 中，Drinfeld 对 $G=\GL_{2}$ 的 Drinfeld 例 \ref{ex:Dr} 构造了 $\Sht^{(I_{1},I_{2}),(\l_{1},\l_{2})}_{\GL_{2}, \bK}$（的截断）的紧化。在 \cite{Laff-comp} 中，无水平结构时，Drinfeld 的构造被推广到 $\GL_{n}$。在 \cite{LLaff} 中，带水平结构的 Drinfeld 例 \ref{ex:Dr} 的紧化被构造出来（性质稍逊），并用于证明函数域上的整体 Langlands 对应。在 \cite{NDT} 中，Ng\^o Dac Tuan 用几何不变量理论给出了紧化的另一种构造。

这些构造的基本想法是：放宽 \S\ref{sss:Sht points} 的模描述中的同构 $\io: \cE_{r}\isom {}^{\t}\cE_{0}$，把它换成一个 $G$ 所附的 Vinberg 幺半群中的点。当 $G=\GL_{n}$ 时，这样的点通常称为两个向量空间或向量丛之间的完全同态（complete homomorphism）。对一般的 $G$，Ng\^o Dac Tuan 与 Varshavsky 的一份未发表手稿沿这条路线给出了 shtuka 模空间紧化的构造。

据作者所知，描述 $\Sht^{(I_{1},\cdots, I_{r}),\le\un\l}_{G,\bK}$ 的紧化的奇点仍是一个公开问题。

\subsection{代数圈}
我们简要综述 shtuka 模叠上代数圈的一些例子。为简单起见，下面的讨论中不考虑水平结构。

\ssszh{地平圈}
\S\ref{sss:it Hk} 中给出的迭代 Hecke 叠的定义对任意代数群都说得通，不必要求约化。因此迭代 shtuka 的模叠对 $k$ 上任意代数群也有定义。我们把这一构造用于抛物子群 $P\subset G$，其 Levi 商为 $L$。设有限集 $I$ 有分解 $I=I_{1}\coprod \cdots \coprod I_{r}$，$S$ 是 $k$ 上的测试概形，则 $\Sht^{(I_{1},\cdots, I_{r})}_{P}(S)$ 分类如下的组
$$((x_{i})_{i\in I}, (\cE_{P,j})_{0\le j\le r}, (\a_{j}: \cE_{P, j-1}\dr\cE_{P, j})_{1\le j\le r}, \io)$$
这里与 \S\ref{sss:Sht points} 中 $\Sht^{(I_{1},\cdots, I_{r})}_{G}(S)$ 的描述唯一的差别是：$\cE_{P,j}$ 是 $X$ 上的 $P$-主丛。由该构造对换群的函子性，典范映射 $G\leftarrow{P}\to L$ 诱导映射
\begin{equation*}
\begin{tikzcd}
\Sht_{G}^{(I_{1},\cdots, I_{r})} & \Sht_{P}^{(I_{1},\cdots, I_{r})} \arrow[l,"\frp"'] \arrow[r,"\frq"] & \Sht_{L}^{(I_{1},\cdots, I_{r})}.
\end{tikzcd}
\end{equation*}
映射 $\frp$ 把 $\cE_{P,j}$ 变成诱导的 $G$-主丛 $\cE_{j}=\cE_{P,j}\times^{P}G$；$\frq$ 把 $\cE_{P,j}$ 变成诱导的 $L$-主丛 $\cE_{L,j}=\cE_{P,j}/U_{P}$（其中 $U_{P}$ 是 $P$ 的幂幺根）。

要定义 $\Sht^{(I_{1},\cdots, I_{r})}_{P}$ 的有界版本，只需取 $L$ 的一族容许的支配余权 $\un\mu=(\mu_{i})_{i\in I}$，并令
\begin{equation*}
\Sht_{P}^{(I_{1},\cdots, I_{r}), \le\un\mu}=\frq^{-1}(\Sht_{L}^{(I_{1},\cdots, I_{r}), \le\un\mu} ).
\end{equation*}
\cite[Proposition 5.7]{Var} 证明了：若限制到 $U^{I}$ 的\emph{一般点 $\y^{I}$} 之上，则映射
$$\frp^{\le\un\mu}_{\y^{I}}: \Sht_{P,\y^{I}}^{(I_{1},\cdots, I_{r}), \le\un\mu}\to \Sht_{G,\y^{I}}^{(I_{1},\cdots, I_{r})}$$
（即 $\frp$ 的限制）是有限且非分歧的。$\Sht_{P,\y^{I}}^{(I_{1},\cdots, I_{r}), \le\un\mu}$ 的基本圈的像定义了 $\Sht_{G,\y^{I}}^{(I_{1},\cdots, I_{r})}$ 上的\emph{地平圈}。

粗略地说，地平圈的圈类对应 $\Sht_{G}^{(I_{1},\cdots, I_{r})}$ 上同调中的 Eisenstein 部分。对 \ref{ex:Dr} 中的 $\GL_{2}$ 例子，地平圈的几何在 \cite[\S4]{Dr-FBun} 中被分析，它与 Eisenstein 级数的关系在 \cite[\S3]{Dr-Pet} 中有详细研究。

一般而言，按自守形式空间谱分解的风格来描述 $\Sht_{G}^{(I_{1},\cdots, I_{r}),\le\un\l}$ 上同调的 Eisenstein 部分，仍是一个公开问题。$G=\GL_{2}$ 情形的定性陈述见 \cite[\S7]{YZ1}，一般 $G$ 的有限性结果见 \cite{X}。

\ssszh{特殊圈}

Shimura 簇上的特殊圈由 Shimura 子簇给出，也就是说，它们来自 $G$ 的约化子群。类似地可以讨论 shtuka 模叠上的特殊圈。具体地说，若 $\ph: H\to G$ 是 $k$ 上约化群之间的同态，则由 shtuka 概念对换群的函子性，它诱导模叠之间的映射
$$\th: \Sht^{(I_{1},\cdots, I_{r})}_{H}\to \Sht^{(I_{1},\cdots, I_{r})}_{G}.$$

对有界版本，可以先取 $H$ 的一族支配余权 $\un\mu$，再为 $G$ 取足够大的 $\un\l$，使 $\th$ 把 $\Sht^{(I_{1},\cdots, I_{r}), \le\un\nu}_{H}$ 送到 $\Sht^{(I_{1},\cdots, I_{r}), \le\un\l}_{G}$。Breutmann \cite{Br} 证明了：在允许对 $\Sht^{(I_{1},\cdots, I_{r})}_{H}$ 加抛孔水平结构的一般情形下，$\th$ 是概形的、有限且非分歧的。另见 \cite{Y-cycle} 中的另一个证明。于是对适当的 $\un\l$，$\th$ 的像定义了 $\Sht^{(I_{1},\cdots, I_{r}),\le\un\l}_{G}$ 的 Chow 群中的一个元素。

我们给出几个具有算术意义的特殊圈的例。其中有些并不严格落在上述情形之内。

\begin{exam}[Heegner–Drinfeld 圈]\label{ex:HD}
设 $\nu: X'\to X$ 是非分歧二重覆叠。设 $r$ 为偶数，并把 $I=\{1,2,\cdots, r\}$ 分解为单点集 $I_{i}=\{i\}$。设 $\un\mu=(\mu_{i})_{1\le i\le r}$，其中每个 $\mu_{i}$ 或是 $(1,0,\cdots,0)$，或是 $(0,\cdots, 0,-1)$（$\GL_{n}$ 的极小余权）。令 $\Sht^{\un\mu}_{n,X'}$ 表示 $X'$ 上秩 $n$ 的迭代 shtuka、界由 $\un\mu$ 给出的模叠 $\Sht^{(I_{1},\cdots, I_{r}),\le\un\mu}_{\GL_{n},X'}$。

设 $\un\l=(\l_{i})_{1\le i\le r}$ 是 $\GL_{2n}$ 的由 $\un\mu$ 决定的极小支配余权组：若 $\mu_{i}=(\underbrace{0,\cdots, 0}_{n-1},-1)$，则 $\l_{i}=(\underbrace{0,\cdots,0}_{2n-1}, -1)$；若 $\mu_{i}=(1,\underbrace{0,\cdots,0}_{n-1})$，则 $\l_{i}=(1, \underbrace{0,\cdots,0}_{2n-1})$。类似地，用 $\Sht^{\un\l}_{2n,X}$ 记 $X$ 上秩 $2n$ 的迭代 shtuka、界由 $\un\l$ 给出的模叠 $\Sht^{(I_{1},\cdots, I_{r}),\le\un\l}_{\GL_{2n},X}$。

有映射
\begin{equation*}
\th^{\un\mu}_{\HD(n)}: \Sht^{\un\mu}_{n, X'}\to \Sht^{\un\l}_{2n,X}
\end{equation*}
它把 $X'$ 上秩 $n$ 的 shtuka $((x'_{i}), (\cE_{i}), (\a_{i}), \io)$ 送到 $X$ 上秩 $2n$ 的 shtuka $((\nu(x'_{i})), (\nu_{*}\cE_{i}),(\nu_{*}\a_{i}), \nu_{*}\io)$。把基本圈在 $\th^{\un\mu}_{\HD(n)}$ 下的像称为\emph{Heegner–Drinfeld 圈}。由命题 \ref{p:local model}(2)，$\dim \Sht^{\un\mu}_{n,X'}=rn$，而 $\dim \Sht^{\un\mu}_{2n,X}=2rn$，故有
\begin{equation*}
\dim  \Sht^{\un\mu}_{n,X'}=\frac{1}{2}\dim \Sht^{\un\l}_{2n,X}.
\end{equation*}
因此 $\Sht^{\un\mu}_{n,X'}$ 的基本类在 $\th^{\un\mu}_{\HD(n)}$ 下的像定义了 $\Sht^{\un\l}_{2n,X}$ 上的一个中间维数的圈。

为了看出这个圈与 Heegner 点的关系，我们考虑 $n=1$ 的情形以及上述构造带一条腿的变体。设 $\infty\in |X|$ 是在 $X'$ 中惰性的闭点，于是 $\infty'\in |X'|$ 是 $\infty$ 上方唯一的闭点。考虑 \S\ref{sss:one leg} 中的 Drinfeld 模（相对）曲线 $\Sht^{\Dr}_{2}$。设 $\un\mu=(\mu_{1},\mu_{2})$，其中作为 $\Gm$ 的余权 $\mu_{1}=-1$，$\mu_{2}=1$。令 $\cP_{X'}=\Sht^{(\{1\},\{2\}), \un\mu}_{\GL_{1}, X'}|_{U'\times \{\infty'\}}$；这是 $U'=X'-\infty'$ 上的 $\Pic_{X'}(k)$-主丛。我们构造映射
\begin{equation*}
\th^{1}_{\HD(1)}: \cP_{X'}\to \Sht^{\Dr}_{2}.
\end{equation*}
对 $\cP_{X'}$ 的一个 $S$-点 $\cL_{0}\leftarrow\cL_{0}(-x')\rightarrow\cL_{0}(-x'+y')\isom {}^{\t}\cL_{0}$（其中 $y':S\to \infty'$），考虑 $X'$ 上秩二的丛 $\cE=\nu_{*}\cL_{0}$，它在 $S\times\infty$ 上的纤维有分裂 $\cE|_{S\times\infty}=\cE_{y'}\op \cE_{\s(y')}$（其中 $\s\in \Gal(X'/X)$ 是非平凡元素）。给 $\cE|_{S\times \infty}$ 配上完全旗 $\cE_{y'}\subset \cE_{S\times\infty}$，并定义 $\cE(-\frac{1}{2})\subset \cE$ 为 $\cE_{y'}$ 在限制映射 $\cE\to \cE|_{S\times \infty}$ 下的原像。于是 $\cE(-\frac{1}{2})=\nu_{*}\cL_{0}(-\s(y'))$。因此 ${}^{\t}(\cE(-\frac{1}{2}))=\nu_{*}({}^{\t}\cL_{0})(-y')\cong \nu_{*}\cL_{0}(-x')$（这里用到 $y'\c\Fr_{S}=\s(y')$，因为 $\infty$ 是惰性的）。设 $\phi: {}^{\t}(\cE(-\frac{1}{2}))\to \cE$ 由嵌入 $\cL_{0}(-x')\incl \cL_{0}$ 诱导。这就定义了 $S$-点 $(\xi=\nu(x'), \cE, \{\cE(\frac{i}{2})\}, \phi)\in \Sht^{\Dr}_{2}(S)$。

利用 \S\ref{sss:dict} 中椭圆模与 shtuka 之间的词典，$\cP_{X'}$ 分类 $X$ 上秩二的椭圆模，其自同态由 $A'=\G(U', \cO_{U'})$ 给出；这是带复乘的椭圆曲线的函数域类比。因此 $\th(\cP_{X'})$ 是模曲线上 Heegner 点的函数域类比。带 Iwahori 水平结构的推广在 \cite{YZ2} 中完成。
\end{exam}

\begin{exam}[Gan–Gross–Prasad 圈]\label{ex:GGP}
设 $\nu: X'\to X$ 是非分歧二重覆叠。设 $\Ug_{n}$ 是 $X$ 上关于 $X'$ 的酉群概形。我们不细说非分裂群 $\Ug_{n}$ 本身，只把 $\Bun_{\Ug_{n}}$ 与 $\Sht^{r}_{\Ug_{n}}$ 的模含义写清楚。

设 $S$ 是 $k$ 上的测试概形。$X \times S$ 上关于 $\nu:X'\to X$ 的秩 $n$ \emph{Hermite}（Hermitian，亦称\emph{酉}）丛，是指 $X' \times S$ 上秩 $n$ 的向量丛 $\cF$，并带有同构 $h : \cF \xrightarrow{\sim} \sigma^* \cF^{\vee}$，满足 $\sigma^* h^{\vee} = h$。这里 $\cF^{\vee}=\un\Hom(\cF, \om_{X'})$ 是 $\cF$ 的 Serre 对偶。于是 $\Bun_{\Ug_{n}}$ 是这样的叠：其 $S$-点是 $X \times S$ 上关于 $\nu$ 的秩 $n$ Hermite 丛构成的群胚。

$X'\times S$ 上秩 $n$ Hermite 丛 $(\cF,h)$ 沿 $x': S\to X'$ 的一个初等下修正，是指另一个秩 $n$ Hermite 丛 $(\cF',h')$（同在 $X'\times S$ 上），连同同构 $\a:  \cF|_{X' \times S - \Gamma(x') - \Gamma(\sigma(x'))} \xrightarrow{\sim} \cF'|_{X' \times S - \Gamma(x')-\Gamma(\sigma(x'))}$，它与 Hermite 结构相容，并满足下述性质：存在秩 $n$ 的向量丛 $\cF^{\flat}$（必然唯一），使 $f$ 可以分解为
\begin{equation*}
\begin{tikzcd}
\cF & \cF^{\flat} \arrow[l, hook, "\oll{\a}"] \arrow[r, hook, "\orr{\a}"] & \cF'
\end{tikzcd}
\end{equation*}
使得 $\coker(\oll{\a})$ 在 $\Gamma(x')$ 上是秩 $1$ 局部自由，$\coker(\orr{\a})$ 在 $\Gamma(\sigma(x'))$ 上是秩 $1$ 局部自由。

设 $r$ 为偶数。定义叠 $\Sht_{\Ug_{n}}^{r}$，其 $S$-点由下述数据构成的群胚给出
$$((x'_{i})_{1\le i\le r}, (\cF_{i}, h_{i})_{0\le i\le r}, (\a_{i})_{1\le i\le r}, \io)$$
其中 $x'_{i}: S\to X'$，$(\cF_{i},h_{i})$ 是 $X'\times S$ 上秩 $n$ 的 Hermite 丛，$\a_{i}: \cF_{i-1}\dr\cF_{i}$ 是 $\cF_{i-1}$ 沿 $x'_{i}$ 的初等下修正，$\io$ 是 Hermite 丛的同构 $\cF_{r}\cong {}^{\t}\cF_{0}$。

现在固定 $X'$ 上关于 $\nu$ 的秩一 Hermite 丛 $\cL$；这等价于给一个线丛以及同构 $\Nm_{X'/X}\cL\cong \om_{X}$。有映射
\begin{equation*}
\ka_{\cL}:\Sht^{r}_{\Ug_{n}}\to \Sht^{r}_{\Ug_{n+1}}
\end{equation*}
它把 $((x'_{i}), (\cF_{i}), \cdots)$ 送到 $((x'_{i}), (\cF_{i}\op\cL),\cdots)$。考虑对角映射
\begin{equation*}
\th^{r}_{\GGP(n)}=(\id,\ka_{\cL}): \Sht^{r}_{\Ug_{n}}\to \Sht^{r}_{\Ug_{n}}\times_{X'^{r}}\Sht^{r}_{\Ug_{n+1}}.
\end{equation*}
由 \cite[Lemma 6.9(2)]{FYZ1}，映射 $\Sht^{r}_{\Ug_{n}}\to X'^{r}$ 光滑，相对维数为 $r(n-1)$。因此
$$\dim\Sht^{r}_{\Ug_{n}}=\frac{1}{2}(\dim\Sht^{r}_{\Ug_{n}}\times_{X'^{r}}\Sht^{r}_{\Ug_{n+1}}).$$
$\Sht^{r}_{\Ug_{n}}$ 的基本圈在 $\th^{r}_{\GGP(n)}$ 下的像定义了目标上的一个中间维数的圈。记其圈类为 $Z^{r}_{\GGP(n)}\in \cohog{2rn}{(\Sht^{r}_{\Ug_{n}}\times_{X'^{r}}\Sht^{r}_{\Ug_{n+1}})_{\ov k}}(rn)$。

受算术 Gan–Gross–Prasad 猜想 \cite{GGP} 的启发，我们预期：对 $\Ug_{n}(\AA_{F})\times \Ug_{n+1}(\AA_{F})$ 的尖点自守表示 $\Pi$，$Z^{r}_{\GGP(n)}$ 的 $\Pi$-等型分支的自交数可以良定义（注意它先验地并不是具有 proper 支集的圈），并且与 $\Pi$ 到 $\GL_{n}(\AA_{F'})\times\GL_{n+1}(\AA_{F'})$ 的基变换的 Rankin–Selberg $L$-函数的第 $r$ 个中心导数有关，其中 $F'=k(X')$。

\end{exam}

\begin{exam}[Kudla–Rapoport 圈]\label{ex:KR}
在 \cite{FYZ1} 中我们引入了一类代数圈，它们是酉 Shimura 簇上 Kudla–Rapoport 特殊圈 \cite{KR2} 的函数域类比。

沿用例 \ref{ex:GGP} 的设定。设 $\cE$ 是 $X'$ 上秩 $m$ 的向量丛。定义叠 $\cZ_{\KR(n)}^{r}(\cE)$，其 $S$-点是下述数据构成的群胚：
 \begin{itemize}
 \item 一个 Hermite shtuka $((x'_i)_{1\le i\le r}, (\cF_{i},h_{i})_{0\le i\le r}, (\a_{i})_{1\le i\le r}, \io) \in  \Sht_{\Ug_{n}}^{r}(S)$。
 \item $X'\times S$ 上凝聚层之间的映射 $t_i : \cE \boxtimes \cO_{S} \to \cF_i$（$i=0,1,\cdots, r$），使得同构 $\io : \cF_r \cong {}^{\t}\cF_0$ 把 $t_r$ 与 ${}^{\t}t_0$ 相缠，且对每个 $i = 1, \cdots, r$，映射 $t_{i-1}, t_{i}$ 被修正 $\a_i :  \cF_{i-1} \dr \cF_{i}$ 相缠，即下图交换。
\begin{equation}\label{EFsq}
\begin{tikzcd}
\cE\boxtimes\cO_{S} \arrow[r,equal] \arrow[d,"t_{0}"'] & \cE\boxtimes\cO_{S} \arrow[r,equal] \arrow[d,"t_1"'] & \cdots \arrow[r,equal] \arrow[d] & \cE\boxtimes\cO_{S} \arrow[r,"\sim"] \arrow[d,"t_{r}"'] & {}^{\t}(\cE \boxtimes\cO_{S}) \arrow[d,"{}^{\t}t_{0}"] \\
\cF_0 \arrow[r,dashed,"\a_0"'] & \cF_1 \arrow[r,dashed,"\a_1"'] & \cdots \arrow[r,dashed,"\a_r"'] & \cF_{r} \arrow[r,"\io","\sim"'] & {}^{\t}\cF_0
\end{tikzcd}
\end{equation}
 \end{itemize}

对 $((x'_i), (\cF_i,h_{i}), (\a_i), \io, (t_i)) \in \cZ_{\KR(n)}^r(\cE)(S)$，考虑复合
\begin{equation*}
\cE \boxtimes \cO_S \xr{t_i} \cF_i \xr{h_i} \s^* \cF_i^{\vee} \xr{\s^* t^{\vee}_{i}} \s^* \cE^{\vee} \boxtimes \cO_S.
\end{equation*}
由 \eqref{EFsq} 中的交换方块，上述映射与 $i$ 无关，并且下降为映射 $a: \cE\to \s^{*}\cE^{\vee}$，满足 $\s^{*}a^{\vee}=a$。换言之，$\cZ^{r}_{\KR(n)}(\cE)$ 的每个点都在 $\cE$ 上诱导一个可能退化的 Hermite 形式；据此有开闭子叠的分解
\begin{equation*}
\cZ_{\KR(n)}^{r}(\cE)=\coprod_{a\in \Herm(\cE)} \cZ^{r}_{\KR(n)}(\cE,a).
\end{equation*}
其中 $\Herm(\cE)$ 是 $\cE$ 上（可能退化的）Hermite 形式之集。有自然映射
\begin{equation*}
\th_{\KR(n)}^{r}(\cE,a): \cZ^{r}_{\KR(n)}(\cE,a)\to \Sht^{r}_{\Ug_{n}}
\end{equation*}
它是有限且非分歧的。

$\cZ^{r}_{\KR(n)}(\cE,a)$ 的预期维数是 $r(n-m)$。然而它的实际维数往往更大。在 \cite{FYZ2} 中，我们用导出代数几何给 $\cZ^{r}_{\KR(n)}(\cE,a)$ 配上一个预期维数 $r(n-m)$ 的虚拟基本类 $[\cZ^{r}_{\KR(n)}(\cE,a)]$；它在 $\th^{r}_{\KR(n)}(\cE,a)$ 下的像是酉 Shimura 簇上 Kudla–Rapoport 特殊圈的函数域类比。
\end{exam}

\subsection{相交理论与 $L$-函数的高阶导数}
经典上，自守形式的周期积分常常可以表示成自守 $L$-函数的特殊值。在几何情形中，Shimura 簇或 shtuka 模叠上的特殊圈可以看作周期积分的类比。在数域上，Shimura 簇上特殊圈的高度有时与自守 $L$-函数的导数有关，例如在 Gross–Zagier 公式或算术 Gan–Gross–Prasad 猜想中。在函数域上，不仅有类似的故事，还有进一步的推广：shtuka 可以有多条腿，这种灵活性使人能把特殊圈的相交数与 $L$-函数的高阶导数联系起来。我们简要综述这一方向的若干结果。关于动机与证明中所用技术的更详细讨论，参见 \cite{Y-ICM}。

\ssszh{高阶 Gross–Zagier 公式}

著名的 Gross–Zagier 公式用模形式的自守 $L$-函数的一阶导数，给出模曲线上（投影到 Hecke 本征分量的）Heegner 点高度之间的等式。该公式在数论中有大量应用，尤其在椭圆曲线的算术中。

在 \cite{YZ1} 与 \cite{YZ2} 中，我们给出 Gross–Zagier 公式的函数域类比的高阶导数／多条腿版本。先陈述无水平结构的版本，它在 \cite{YZ1} 中处理。

设 $G=\PGL_{2}$，$I=\{1,\cdots, r\}$ 分解为单点集 $I_{i}=\{i\}$，并设 $\l_{i}$ 是 $G$ 唯一的支配极小余权。把所得的 shtuka 模叠 $\Sht^{(I_{1},\cdots, I_{r}),\le\un\l}_{G}$ 简记为 $\Sht^{r}_{G}$。$\un\l$ 的容许性条件迫使 $r$ 为偶数。

设 $\nu:X'\to X$ 是非分歧二重覆叠。设 $T=(\Res_{X'/X}(\Gm\times X'))/\Gm$ 是 $X$ 上在 $X'$ 处分裂的一维非分裂环面。对 $X$ 上这样的群概形 $T$ 与序列 $\un\mu=(\mu_{1},\cdots, \mu_{r})\in\{\pm1\}^{r}$，有定义 $\Sht_{T}^{\un\mu}:=\Sht_{T}^{(I_{1},\cdots, I_{r}), \un\mu}$。具体地说，它可以等同于 $\Sht^{(I_{1},\cdots, I_{r}), \un\mu}_{\GL_{1}, X'}$（基曲线取 $X'$，见例 \ref{ex:Gm}）在 $\Pic_{X}(k)$ 作用下的商；该作用为拉回到 $X'$ 再与秩一 shtuka 作张量积。如例 \ref{ex:HD} 那样，有自然映射
\begin{equation*}
\th: \Sht^{\un\mu}_{T}\to \Sht^{r}_{G}.
\end{equation*}
考虑 $\th$ 的提升
\begin{equation*}
\th': \Sht^{\un\mu}_{T} \to \Sht'^{r}_{G}:=\Sht^{r}_{G}\times_{X^{r}}X'^{r}.
\end{equation*}
容易看出 $\Sht^{\un\mu}_{T}$ 是 proper 的、维数为 $r$，而由推论 \ref{c:Sht sm}，$\Sht'^{r}_{G}$ 是维数 $2r$ 的光滑概形。因此 Heegner–Drinfeld 圈 $\cZ^{\un\mu}_{T}:=\th'_{*}[\Sht^{\un\mu}_{T}]$ 是光滑 Deligne–Mumford 叠 $\Sht'^{r}_{G}$ 中的中间维数 proper 圈。设 $Z^{\un\mu}_{T}\in \cohoc{2r}{\Sht'^{r}_{G,\kbar},\Qlbar}(r)$ 是 $\cZ^{\un\mu}_{T}$ 的圈类。

现在设 $\pi$ 是 $G(\AA)$ 的处处非分歧尖点自守表示，系数在 $\Qlbar$ 中。在 Hecke 代数作用下，$Z^{\un\mu}_{T}$ 到 $\pi$-等型直和项上的投影是有意义的。把所得的上同调类记作 $Z^{\un\mu}_{T,\pi}\in \cohoc{2r}{\Sht'^{r}_{G,\kbar},\Qlbar}(r)$。

\begin{thm}[\cite{YZ1}]\label{th:YZ1} 圈类 $Z^{\un\mu}_{T,\pi}$ 的自交数为
\begin{equation*}
\j{Z^{\un\mu}_{T,\pi}, Z^{\un\mu}_{T,\pi}}_{\Sht'^{r}_{G}}=\frac{q^{2-2g}}{2(\log q)^{r}}\frac{\sL^{(r)}(\pi_{F'}, 1/2)}{L(\pi,\Ad,1)}
\end{equation*}
其中
\begin{itemize}
\item $\pi_{F'}$ 是 $\pi$ 到 $F'=k(X')$ 的基变换。
\item $\sL(\pi_{F'},s)=q^{4(g-1)(s-1/2)}L(\pi_{F'},s)$ 是 $\pi_{F'}$ 的规范化 $L$-函数，满足 $\sL(\pi_{F'},s)=\sL(\pi_{F'},1-s)$。
\end{itemize}
\end{thm}

在 \cite{YZ2} 中，我们把上述定理推广到允许自守表示 $\pi$ 带 Iwahori 水平结构（即局部表示 $\pi_{v}$ 或是非分歧的，或是 Steinberg 表示的非分歧扭），也允许二重覆叠 $\nu:X'\to X$ 有分歧。这一般性使将来能把结果应用于函数域上的 Birch–Swinnerton-Dyer 猜想，预览见 \cite[\S1.3]{YZ2}。

\ssszh{高阶 Siegel–Weil 公式与高阶 theta 级数}
经典的 Siegel–Weil 公式把辛群（相应地，酉群）上 Siegel–Eisenstein 级数的特殊值与 theta 函数联系起来；后者是数域上二次型（相应地，Hermite 型）表示数的生成级数。在 \cite{Kud} 中，Kudla 开始研究 Siegel–Weil 公式的算术版本：他定义了「算术 theta 函数」——Shimura 曲线整模型上算术圈的生成级数——并发现它与 $\Sp_{4}$ 上 Siegel–Eisenstein 级数的第一个中心导数的关系。在一系列论文 \cite{KR1} 与 \cite{KR2} 中，Kudla 与 Rapoport 通过定义正交群与酉群的 Shimura 簇整模型上特殊圈（Kudla–Rapoport 圈）的生成级数，发展了这一范式。他们猜想这些圈的算术相交数与 Siegel–Eisenstein 级数中心导数的非奇异 Fourier 系数之间的关系。这一猜想最近由 Chao Li 与 Wei Zhang 在 \cite{LZ} 中证明。

在 \cite{FYZ1} 中，我们证明了函数域上酉群的 Siegel–Weil 公式的高阶导数版本。考虑例 \ref{ex:KR} 中构造的特殊圈 $\cZ^{r}_{\KR(n)}(\cE,a)$。我们限定在 $\rk (\cE)=n$ 且 $a: \cE\to\s^{*}\cE^{\vee}$ 在一般点上是同构的情形（这样的 $a$ 称为非奇异的）。此时 $\cZ^{r}_{\KR(n)}(\cE,a)$ 是 $k$ 上的 proper 概形，但它的维数可能不等于预期维数（此时预期维数为 $0$）。我们先定义虚拟基本圈 $[\cZ^{r}_{\KR(n)}(\cE,a)]$，它是 proper 概形 $\cZ^{r}_{\KR(n)}(\cE,a)$ 上的 $0$-圈，从而其次数有定义。我们证明下述高阶 Siegel–Weil 公式

\begin{thm} 设 $r\in \NN$ 为偶数，$\cE$ 是 $X'$ 上秩 $n$ 的向量丛，$a: \cE\to \s^{*}\cE^{\vee}$ 是 $\cE$ 上的非奇异 Hermite 形式。则
\begin{equation*}
\deg [\cZ_{\KR(n)}^r(\cE, a)]=\frac{1}{(\log q)^r}  \left(\dfrac{d}{ds} \right)^r\Big|_{s=0} \left( q^{ds} \wt{E}_{a}(\cE,s,\Phi) \right).
\end{equation*}
这里 $d=-\deg(\cE)+n\deg\omega_X =-\chi(X',\cE)$，而 $\wt{E}_{a}(\cE,s,\Phi)$ 是酉群 $\Ug_{2n}$ 上 Siegel–Eisenstein 级数 $E(g,s,\Phi)$ 的第 $a$ 个 Fourier 系数的某种规范化，在 $\cE$ 处展开（把 $\cE$ 提升为 Siegel–Levi $\GL_{n}(\AA_{F})$ 中的元素）。
\end{thm}

证明与数域情形完全不同（后者基于对 $n$ 的归纳）。证明化归为一个来自 Springer 理论的反常层等式。

在 \cite{FYZ2} 中，我们构造特殊圈的完整生成级数，即对任意秩 $m\le n$ 的 $\cE$ 与 $\cE$ 上任意的 Hermite 形式 $a$，定义虚拟基本类 $[\cZ_{\KR(n)}^r(\cE, a)]$。特别地，$a$ 可以是奇异的甚至为零，此时虚拟基本类会涉及 $\Sht^{r}_{\Ug_{n}}$ 上万有丛的 Chern 类。令 $Z_{\KR(n)}^{r}(\cE,a)$ 为 $[\cZ_{\KR(n)}^r(\cE, a)]$ 在映射 $\th^{r}_{\KR(n)}(\cE,a)$ 下的像，看作 Chow 群 $\Ch_{r(n-m)}(\Sht^{r}_{\Ug_{n}})$ 中的元素。然后把这些代数圈装配成高阶 theta 级数。

更确切地说，固定 $m\le n$，考虑 $X$ 上关于二重覆叠 $X'/X$ 的拟分裂酉群 $\Ug_{2m}$ 及其标准的 Siegel 抛物子群 $P_m$。定义函数
\begin{equation*}
\wt Z^{r}_{m}: \Bun_{P_{m}}(k)\to \Ch_{r(n-m)}(\Sht^{r}_{\Ug_{n}})
\end{equation*}
如下。$\Bun_{P_{m}}(k)$ 中的一个点等同于一对 $(\cG,\cE)$，其中 $\cG$ 是 $X'$ 上秩 $2m$ 的 Hermite 向量丛（这里 Hermite 形式指的是同构 $\cG\isom \s^{*}\cG^{*}$，其中 $\cG^{*}$ 是 $\cG$ 的\emph{线性对偶}，而非 Serre 对偶），$\cE$ 是 $\cG$ 的 Lagrangian 子丛。然后用「Fourier 展开」定义 $\wt Z^{r}_{m}(\cG,\cE)$：
\begin{equation*}
\wt Z^{r}_{m}(\cG,\cE) =\chi(\det\cE)q^{n(\deg \cE-\deg\om_{X}) /2}\sum_{a\in \Herm(\cE)}\psi_{0}(\j{e_{\cG,\cE}, a})Z^{ r}_{\KR(n)}(\cE, a).
\end{equation*}
这里 $\psi_{0}:k\to \Qlbar^{\times}$ 是非平凡特征，$e_{\cG,\cE}\in \Ext^{1}(\s^{*}\cE^{*}, \cE)$ 是正合列 $\cE\incl \cG\surj \s^{*}\cE$ 的扩张类；在 Serre 对偶下它可以与 $a\in \Hom(\cE,\s^{*}\cE^{\vee})$ 配对。

\begin{conj}[模性猜想，{\cite[Conjecture 4.12]{FYZ2}}] 函数 $\wt Z^{r}_{m}$ 下降为函数
\begin{equation*}
Z^{r}_{m}:\Bun_{\Ug_{2m}}(k)\to \Ch_{r(n-m)}(\Sht^{r}_{\Ug_{n}}).
\end{equation*}
\end{conj}
换言之，类 $\wt Z^{r}_{m}(\cG,\cE)\in \Ch_{r(n-m)}(\Sht^{r}_{\Ug_{n}})$ 应当只依赖 Hermite 丛 $\cG$，而不依赖它的 Lagrangian 子丛 $\cE$。当 $r=0$ 时，$\Ch_{0}(\Sht^{0}_{\Ug_{n}})$ 就是 $\Bun_{\Ug_{n}}(k)$ 上的函数空间，经典的 Weil 表示理论表明此时 $\wt Z^{r}_{m}$ 确实下降到 $\Bun_{\Ug_{2m}}(k)$。所得 $\Bun_{\Ug_{2m}}(k)\times \Bun_{\Ug_{n}}(k)$ 上的二元函数就是经典的 theta 函数。

在 \cite{FYZ2} 中我们给出了模性猜想的若干证据；在 \cite{FYZ3} 中我们证明：若限制到 $\Sht^{r}_{\Ug_{n}}\to X'^{r}$ 的通用纤维上，并过渡到圈的上同调类，则猜想成立。
